% annexes/execution.tex
% Exemple d'execution du projet - Analyse du reseau WETH Polygon

\section{Exemple d'ex\'ecution du projet}
\label{ann:execution}

Cette annexe pr\'esente les commandes n\'ecessaires pour reproduire l'ensemble des analyses et compiler le rapport, accompagn\'ees des sorties console attendues.

% ---- Installation des dependances ----

\subsection{Installation des d\'ependances}

L'installation des biblioth\`eques Python s'effectue via \texttt{pip}~:

\begin{lstlisting}[language={}, numbers=none, caption={Installation des d\'ependances Python}]
$ pip install networkx~=3.4.2 matplotlib~=3.9.4 pandas~=2.2.3 \
    numpy~=2.1 scipy~=1.14 powerlaw~=2.0.0
\end{lstlisting}

Un fichier \texttt{requirements.txt} est \'egalement fourni \`a la racine du projet pour simplifier l'installation~:

\begin{lstlisting}[language={}, numbers=none, caption={Installation via requirements.txt}]
$ pip install -r requirements.txt
\end{lstlisting}

% ---- Execution des analyses ----

\subsection{Ex\'ecution des analyses}

Le script orchestrateur \texttt{run\_all.py} lance s\'equentiellement les six analyses de questions de recherche ainsi que les cinq analyses compl\'ementaires (annexe \S\ref{sec:complementaires}). Le fichier CSV \texttt{polygon\_weth\_crash\_2024-08-05.csv} (74~Mo) doit \^etre pr\'esent \`a la racine du projet.

\begin{lstlisting}[language={}, numbers=none, caption={Lancement de l'orchestrateur}]
$ python run_all.py
\end{lstlisting}

La sortie console attendue est la suivante~:

\begin{lstlisting}[language={}, numbers=none, caption={Sortie console de run\_all.py (extraits, ex\'ecution r\'eelle)}]
============================================================
  Analyse du reseau WETH Polygon - Lundi Noir (5 aout 2024)
============================================================

Chargement des donnees...
Graphe construit: 12880 noeuds, 39999 aretes
Version non-orientee: 33827 aretes

============================================================
  RQ1 - Topologie
============================================================
  Noeuds:      12 880      Aretes:      39 999
  Densite:     0.000241    Degre moyen: 3.11
  Composantes WCC: 95   Geante: 12 538 (97.3%)
  Composantes SCC: 7 632 Geante:  5 206 (40.4%)
  Diametre (SCC):           15
  Longueur moyenne chemins: 3.970
Figure sauvegardee: rq1_degree_distribution.png
Figure sauvegardee: rq1_components.png
Figure sauvegardee: rq1_hub_network.png

============================================================
  RQ2 - Centralite
============================================================
  Top degre entrant  : 0xe50fa9...8128c8  (0.0672)
  Top degre sortant  : 0x281805...cb7116  (0.1555)
  Top betweenness    : 0x281805...cb7116  (0.0816)
  Top closeness      : 0x86f1d8...e618e0  (0.4287)
  Top PageRank       : 0xe50fa9...8128c8  (0.0425)
  Spearman degre/PR:        0.6061
  Spearman degre/betw:      0.5144
  Spearman closeness/PR:    0.0080
Figures: rq2_centrality_{barplots,comparison,network}.png

============================================================
  RQ3 - Communautes
============================================================
  Louvain: 174 communautes detectees
  Top 5 tailles: 2981, 2414, 2338, 1542, 1381 noeuds
  Modularite Q = 0.6445  (structure significative)
  Singletons: 3   Communautes >= 10: 26
Figures: rq3_{communities,community_sizes}.png

============================================================
  RQ4 - Temporel
============================================================
  24 fenetres horaires UTC analysees
  Pic transactions: 13h UTC  (28 146 tx)
  Pic volume WETH:  01h UTC  (21 517.99 WETH)
  Pic aretes:       01h UTC  (5 791 aretes)
Figures: rq4_{temporal_evolution,heatmap}.png

============================================================
  RQ5 - Petit-monde
============================================================
  Transitivite (C):  0.029715
  C aleatoire (ER):  0.000381
  Longueur moy (L):  3.9700      L aleatoire: 5.9034
  Sigma = (C/Cr)/(L/Lr) = 115.90
  -> Proprietes petit-monde confirmees
Figures: rq5_{small_world_comparison,clustering_distribution,clustering_network}.png

============================================================
  RQ6 - Flux ponderes
============================================================
  Coefficient de Gini: 0.9780
  Top-10 whales:       24.4% du volume total
  Sous-graphe whales:  43 noeuds, 50 aretes
Figures: rq6_{weight_distribution,lorenz_curve,whale_network}.png

============================================================
  Analyses complementaires
============================================================
LCC: 12 538 noeuds, 33 576 aretes

[1] Ponts (Tarjan)
    5 890 ponts detectes (17.54% des aretes)
    Degre moyen aux extremites: 869.06
[2] Decomposition en k-coeurs (Seidman)
    k_max = 37  (111 noeuds dans le 37-coeur)
    2-coeur: 6 652 noeuds (53.1% de la LCC)
[3] Similarite topologique (Jaccard) sur top 50 hubs
    Paire 0x506c64.../0xf08650...:  J = 1.0000  (265 voisins identiques)
    Mediane des J non nuls: 0.1071
[4] Propagation SI depuis le hub principal (degre 2004)
    p = 0.10 : 78.7% de la LCC apres 15 pas
    p = 0.30 : 99.1% de la LCC apres 15 pas
    p = 0.50 : 100.0% de la LCC apres 15 pas
[5] Robustesse par suppression cascade (100 noeuds retires)
    Aleatoire   : LCC perd  0.90 %
    Top degre   : LCC perd 56.67 %
    Top betweenness : LCC perd 58.35 %
Figures: extra_{bridges_network,kcore_distribution,jaccard_topk,si_propagation,robustness_cascade}.png

============================================================
  Termine en 652.3 secondes
============================================================
\end{lstlisting}

\`A l'issue de l'ex\'ecution, les 21 fichiers PNG sont g\'en\'er\'es dans le r\'epertoire \texttt{figures/} \`a une r\'esolution de 200~DPI~: 16 figures pour les six RQ et 5 figures pour les analyses compl\'ementaires (annexe \S\ref{sec:complementaires}). La dur\'ee totale d'ex\'ecution est d'environ 11 minutes sur un poste de bureau standard, la moiti\'e \'etant consomm\'ee par la d\'etection des ponts ($\mathcal{O}(V+E)$ mais sur un graphe de $33\,576$ ar\^etes) et la simulation de robustesse (recalcul de la LCC \`a chaque retrait).

% ---- Compilation du rapport ----

\subsection{Compilation du rapport}

La compilation du rapport LaTeX n\'ecessite trois passes de \texttt{pdflatex} et une passe de \texttt{biber} pour r\'esoudre correctement les r\'ef\'erences crois\'ees et la bibliographie~:

\begin{lstlisting}[language={}, numbers=none, caption={S\'equence de compilation LaTeX}]
$ cd rapport
$ pdflatex main.tex
$ biber main
$ pdflatex main.tex
$ pdflatex main.tex
\end{lstlisting}

Alternativement, l'outil \texttt{latexmk} automatise cette s\'equence en une seule commande~:

\begin{lstlisting}[language={}, numbers=none, caption={Compilation automatis\'ee via latexmk}]
$ cd rapport
$ latexmk -pdf main.tex
\end{lstlisting}

Le fichier \texttt{main.pdf} produit contient l'int\'egralit\'e du rapport avec figures int\'egr\'ees, table des mati\`eres, bibliographie et annexes.
